% page 598 du fac-simile (image p-597 du scan)
% bloc 152.4 x 250.6 mm ; marge gauche 8.0 mm
\pgc{-12.700mm}{0.000mm}{
\l{}
\l{}
\l{}
\l{}
\l{}
\l{}
\l{}
\l{\sou{transepto}. (\sou{arkit.}) Ta parto di kirko exter la navo, qua}
\l{\cel{3}formacas la \textquotedbl{}brakii\textquotedbl{} di la kruco. - DEFR.}
\l{}
\l{\sou{transferar}. (\sou{trans., de,}\cel{1}ad). I. Instalar, establisar,}
\l{\cel{3}de ca loko ad ita. - II. (\sou{yurocienco}). Igar divenar la}
\l{\cel{3}proprietajo di altru. - DEFIS.}
\l{}
\l{\sou{transfigurar}. (\sou{trans.}) Transformar per altrigar la vizajo,}
\l{\cel{3}la traiti extera. -\cel{2}EFIS.}
\l{}
\l{\sou{transfinita}. (\sou{matem.})\cel{1}\sou{Nombri transfinita}\cel{1}: Nombri qui egar-}
\l{\cel{3}desis en lá teorio pri la ensembli. (\sou{Ensemblo matemati}-}
\l{\cel{3}\sou{kala}\cel{1}:\cel{2}Grupo determinita de nombri, objekti, e c.; la}
\l{\cel{3}literi di la alfabeto, la nombri pozitiva ne-para, e c.}
\l{\cel{3}esas ensembli). -}
\l{}
\l{\sou{transformar}. (\sou{trans., ulo, ulu, ad)}. I. Pasigar de ca formo}
\l{\cel{3}ad ita. - II. (\sou{metaf.}) Par-modifikar. - DEFIS.}
\l{}
\l{\sou{transformatoro}. (\sou{tekn.}) (\sou{elektro)}. Aparato qua, recevante}
\l{\cel{3}energio elektrala di ta o ca speco, modifikas exkluzive}
\l{\cel{3}lua formo.- DEFIRS.}
\l{}
\l{\sou{transitar}. (\sou{netrans.}) (\sou{komerco}).(\sou{aludante vari}). Transirar}
\l{\cel{3}lando, urbo, sen pagar (pri la imposti doganala od aciz-}
\l{\cel{3}ala), pro yuro grantita pri lo en ula cirkonstanci).-}
\l{\cel{3}DEFIRS.}
\l{}
\l{\sou{translacar}. (\sou{netrans.}) (\sou{aludante solido}). Movar (o movesar)}
\l{\cel{3}tale ke omna parti havas rapideso-grado e direciono}
\l{\cel{3}konstanta. - EFIS.}
\l{}
\l{\sou{translucida}.\cel{2}Qua lasas lumo pasar tra su, ma tale difuzita}
\l{\cel{3}ke onu ne povas dicernar precize la objekti tra lu.-}
\l{\cel{3}DEFIS.}
\l{}
\l{\sou{transmisar}. (\sou{trans.}) I. (\sou{yurocienco}). Igar pasar ad ulu}
\l{\cel{3}(ta o ca yuro, privilejo, quan onu posedas). - II. (\sou{biol}.}
\l{\cel{3}Igar pasar a la decendanti. - III.(\sou{tekn.})Komunikar la}
\l{\cel{3}movo di organo ad altra organo (per ingranaji, kabli,}
\l{\cel{3}rimeni, kateni, e c.). - DEFIS.}
\l{}
\l{\sou{transmutar}. (\sou{trans., ad)}. (\sou{alkemio}). Chanjar la naturo di}
\l{\cel{3}(metalo ordinara a metalo precoza). - EFIS.}
\l{}
\l{\sou{transparar}.\cel{2}(\sou{netrans}.)\cel{2}(\sou{tra, sub}). Esar videbla tra.}
\l{\cel{3}DEFIS.}
\l{}
\l{\sou{transparento}. Papero-folio lineizita horizontale, quan onu}
\l{\cel{3}lokizas sub la papero-folio sur qua onu skribas, por ke}
\l{\cel{3}la skriburo formacez linei rekta. - DEFIRS.}
}
